\documentclass[../../../include/open-logic-section]{subfiles}

\begin{document}

\olfileid{sth}{ordinals}{wo}
\olsection{خوش‌ترتیب‌ها}

مفهوم بنیادی چنین است:

\begin{defn}
رابطهٔ~$<$، مجموعهٔ~$A$ را \emph{خوش‌ترتیب می‌کند} اگر و تنها اگر این دو
شرط را برآورده سازد:
\begin{enumerate}
	\item $<$ همبند باشد؛ یعنی برای هر~$a,b \in A$، یا~$a<b$ یا~$a=b$
	یا~$b<a$؛
	\item هر زیرمجموعهٔ ناتهی~$A$ دارای یک~!!{element}~$<$-مینیمال باشد؛
	یعنی اگر~$\emptyset \neq X \subseteq A$، آنگاه
	$(\exists m \in X)(\forall z \in X)z \nless m$.
\end{enumerate}
\end{defn}

به‌آسانی می‌توان دید سه نمونه‌ای که اکنون بررسی کردیم واقعاً رابطه‌های
خوش‌ترتیب‌کننده بودند.

\begin{prob}
\Olref[sth][ordinals][idea]{sec} سه ترتیب نمونه روی اعداد طبیعی عرضه کرد.
بررسی کنید که هریک خوش‌ترتیب است.
\end{prob}

در ادامه چند مشاهدهٔ مقدماتی ولی بسیار مهم دربارهٔ خوش‌ترتیبی می‌آید.

\begin{prop}\ollabel{wo:strictorder}
اگر~$<$ مجموعهٔ~$A$ را خوش‌ترتیب کند، آنگاه هر زیرمجموعهٔ ناتهی~$A$ یک
عضو یکتای~$<$-کمینه دارد، و~$<$ ناخودبازتاب، نامتقارن و متعدی است.
\end{prop}

\begin{proof}
اگر~$X$ زیرمجموعه‌ای ناتهی از~$A$ باشد، یک~!!{element}~$<$-مینیمال
$m$ دارد؛ یعنی~$(\forall z \in X)z \nless m$. چون~$<$ همبند است،
$(\forall z \in X)m \leq z$. پس~$m$، ~!!{element}~$<$-کمینهٔ~$X$ است.

برای ناخودبازتابی، $a \in A$ را ثابت بگیرید؛ ~!!{element}~$<$-کمینهٔ~$\{a\}$ خود
$a$ است، پس~$a \nless a$. برای تعدی، اگر~$a<b<c$، از آنجا که
$\{a,b,c\}$ یک~!!{element}~$<$-کمینه دارد، $a<c$. نامتقارنی از
ناخودبازتابی و تعدی نتیجه می‌شود.
\end{proof}

\begin{prop}\ollabel{propwoinduction}
اگر~$<$ مجموعهٔ~$A$ را خوش‌ترتیب کند، آنگاه برای هر فرمول~$\phi(x)$:
\[
	\text{اگر }(\forall a \in A)((\forall b < a)\phi(b) \lif
		\phi(a))\text{، آنگاه }(\forall a \in A)\phi(a).
\]
\end{prop}

\begin{proof}
عکس نقیض را اثبات می‌کنیم. فرض کنید~$\lnot(\forall a \in A)\phi(a)$؛
یعنی~$X = \Setabs{x \in A}{\lnot\phi(x)} \neq \emptyset$. آنگاه~$X$
یک~!!{element}~$<$-مینیمال مانند~$a$ دارد. پس
$(\forall b<a)\phi(b)$، ولی~$\lnot\phi(a)$.
\end{proof}\noindent
این ویژگی اخیر باید اصل استقرای قوی روی اعداد طبیعی را به یادتان آورد؛
یعنی: اگر~$(\forall n \in \omega)((\forall m<n)\phi(m) \lif \phi(n))$،
آنگاه~$(\forall n \in \omega)\phi(n)$. و همین ویژگی خوش‌ترتیبی را به
مفهومی بسیار \emph{استوار} تبدیل می‌کند.\footnote{یادآوری: همهٔ فرمول‌ها
می‌توانند پارامتر داشته باشند (مگر آنکه خلافش به‌صراحت گفته شود).}

\end{document}
